% % Бүлэг 1


% % zuwhun uuriin sedewiin tuhai bichne sedewiin neriig husseneeree udirdana
% \chapter{Хичээлийн ухаалаг тэмдэглэлийн системийн тухай онол, арга зүйн судалгаа} 
% \label{Chapter1} % Энэ бүлэг рүү ишлэл хийх бол \ref{Chapter1} командыг ашигла 

% %-------------------------------------------------------------------------------

% % Агуулгад ашигласан хэвшүүлэлтийн зарим командын тодорхойлолт
% \newcommand{\keyword}[1]{\textbf{#1}}
% \newcommand{\tabhead}[1]{\textbf{#1}}
% \newcommand{\code}[1]{\texttt{#1}}
% \newcommand{\file}[1]{\texttt{\bfseries#1}}
% \newcommand{\option}[1]{\texttt{\itshape#1}}

% %-------------------------------------------------------------------------------
% \section{Хичээлийн ухаалаг тэмдэглэлийн системийн тухай}
% \paragraph{}
% Орчин үеийн боловсролын салбарт мэдээллийн хэмжээ асар хурдацтай өсөж, оюутан сурагчид олон тооны ном, PDF материал, лекцийн тэмдэглэл зэргийг богино хугацаанд уншиж судлах шаардлагатай тулгарч байна. Уламжлалт аргаар тэмдэглэл хөтлөх, шалгалтад бэлдэх нь цаг хугацаа их шаарддагаас гадна сурлагын бүтээмжийг бууруулдаг. Иймд хиймэл оюун ухаан (AI) болон оптик тэмдэгт танигч (OCR) технологид суурилсан ухаалаг системийг хөгжүүлснээр сургалтын материалыг автоматаар боловсруулж, хураангуйлах, өөрийгөө сорих флаш карт (flashcard), асуулт хариулт үүсгэх замаар суралцах үйл явцыг илүү сонирхолтой, үр дүнтэй болгох боломжтой юм. Энэхүү систем нь зөвхөн тэмдэглэл хадгалах бус, суралцагчийн ахиц дэвшлийг хянах, сургалтын ухаалаг туслах байх зорилготой.


% \section{Асуудлын тодорхойлолт}
% \paragraph{}
% Сургалтын үйл явц дахь уламжлалт арга барилд дараах нийтлэг асуудлууд тулгарч байна:
% \begin{itemize}
%     \item \textbf{Мэдээллийн хэт ачаалал:} Оюутнууд олон хуудас урт PDF материал уншиж, гол санааг ялгаж авахад маш их цаг зарцуулдаг. Боловсролын салбарын судалгаанаас үзэхэд оюутнуудын 70 гаруй хувь нь цагийн менежмент болон мэдээллийн хэт их багтаамжаас шалтгаалан өгөгдсөн унших материалыг бүрэн судалж амждаггүй бөгөөд энэ нь танин мэдэхүйн хэт ачаалалд (cognitive overload) хүргэдэг байна \cite{sweller2011cognitive, berry2010student}.
    
%     \item \textbf{Сканнердсан болон зурган материалын бэрхшээл:} Багш нарын өгсөн материал эсвэл ном нь текст хэлбэрээр бус зураг байдлаар байх тохиолдол их бөгөөд үүнээс текст хуулж авах, хайлт хийх боломжгүй байдаг. Мэдээллийн технологийн салбарын тайланд дурдсанаар байгууллага болон боловсролын орчин дахь нийт дижитал мэдээллийн 80 орчим хувь нь бүтэцгүй өгөгдөл буюу зураг, сканнердсан баримт хэлбэрээр оршдог байна \cite{shukla2022unstructured}. Иймд OCR технологи ашиглахгүйгээр эдгээр материалыг бүрэн ашиглах боломж хязгаарлагдмал юм.
    
%     \item \textbf{Өөрийгөө сорих хэрэглэгдэхүүн дутмаг:} Уншсан зүйлээ бататгахын тулд флаш карт болон сорилын асуулт гараар бэлтгэх нь оюутнаас давхар хөдөлмөр шаарддаг. Судлаач Ж.Данлоски (J. Dunlosky)-ын өргөн цар хүрээтэй судалгаагаар "өөрийгөө сорих" болон "завсарлагатай давтамж" нь хамгийн үр дүнтэй сургалтын арга болох нь батлагдсан боловч, оюутнуудын дийлэнх нь цаг хугацаа шаарддаг тул эдгээрээс зайлсхийж, зөвхөн материалыг идэвхгүй дахин унших аргыг сонгодог байна \cite{dunlosky2013improving}.
    
%     \item \textbf{Хяналт ба тогтолцоогүй байдал:} Оюутнууд аль хичээлийн, аль сэдвийг бүрэн ойлгосон, юун дээр алдаж байгаагаа бодитоор хянаж чаддаггүй. Өөрийн суралцах явцыг хянадаг, системтэйгээр төлөвлөдөг оюутнуудын сурлагын амжилт бусдаас 20-30\%-иар өндөр байдаг хэдий ч, уламжлалт тэмдэглэлийн дэвтэр болон энгийн уншигч програмууд нь үүнийг автоматаар хянах боломжийг олгодоггүй \cite{zimmerman2002becoming}.
% \end{itemize}

% Эдгээр асуудлууд нь суралцагчийн цагийг үр дүнгүй үрж, сурах идэвхийг бууруулж байгаа тул AI болон OCR технологид суурилан эдгээрийг автоматжуулж шийдвэрлэх нэгдсэн систем зайлшгүй шаардлагатай байна.
% \section{Оролцогч талууд}
% \paragraph{}
% Дээрх асуудлуудыг шийдвэрлэх системийг ашиглах, түүнд нөлөөлөх оролцогч талуудыг дараах байдлаар тодорхойлов:
% \begin{enumerate}
%     \item \textbf{Оюутан, суралцагч:} Системийн гол хэрэглэгч. Хичээлийн материалаа (текст болон зураг/OCR) оруулж, хураангуйлуулж, өөрийгөө сорьж шалгалт өгөх, чатботтой харилцах замаар сурлагын амжилтаа ахиулах зорилготой.
%     \item \textbf{Багш, мэдлэг хуваалцагч:} Өөрийн бэлтгэсэн ухаалаг хичээл, тэмдэглэлийг "Public" төлөвт шилжүүлж, бусад оюутнуудад линкээр түгээх хүсэлтэй хэрэглэгч.
%     \item \textbf{Системийн админ:} Системийн хэвийн ажиллагааг хангах, хэрэглэгчдийн дата болон AI боловсруулалтын хэрэглээ, статистикийг хянах, удирдах үүрэгтэй этгээд.
% \end{enumerate}


% \section{Хэрэглэгчийн шаардлагын судалгаа}
% \paragraph{}
% Оролцогч талуудын хэрэгцээнд үндэслэн систем дараах үндсэн шаардлагуудыг хангасан байх ёстой. (Эдгээр нь 2-р бүлэгт заасан системийн FR, NFR шаардлагуудын үндэс болно).
% \begin{itemize}
%     \item \textbf{Суралцагчийн шаардлага:} 
%     \begin{itemize}
%         \item Нарийн төвөгтэй бүтэцтэй болон зураг агуулсан (OCR шаардлагатай) файлуудыг саадгүй боловсруулдаг байх.
%         \item AI-аас хичээлийн агуулгын хураангуй, флаш карт, сорилын асуулт гаргуулж авдаг байх.
%         \item Ойлгохгүй зүйлээ тухайн сэдэвтэй холбогдсон хиймэл оюун ухаант туслахаас асуудаг байх.
%         \item Өөрийн суралцах явц буюу ахиц дэвшлийг системээс хянадаг, шалгалтад тэнцсэнээр хичээлийг үзэж дуусгасанд тооцдог байх.
%     \end{itemize}
%     \item \textbf{Багш, контент бүтээгчийн шаардлага:} Өөрийн үүсгэсэн материалыг олон нийтэд нээлттэй болгож хуваалцах боломжтой байх.
%     \item \textbf{Админы шаардлага:} Системийн ачаалал, алдаа болон нийт хэрэглэгчдийн статистикийг харах, хяналтын самбартай байх.
% \end{itemize}

% \section{Ижил төстэй системийн судалгаа}
% \paragraph{} Дэлхийн зах зээл дээр хиймэл оюун ухаанд суурилсан тэмдэглэл хөтлөх болон сургалтын хэд хэдэн системүүд байдаг. Эдгээрийг судалж үзэх нь өөрийн системийн давуу талыг тодорхойлоход чухал юм.

% \begin{enumerate}
%   \item \textbf{Notion AI:} Notion нь тэмдэглэл хөтлөх хүчирхэг хэрэгсэл бөгөөд сүүлийн үед AI нэмэлтийг ашиглан текстийг хураангуйлах, засварлах боломжтой болсон. Гэвч энэ нь сургалтын үйл явцад зориулагдаагүй, ерөнхий зориулалттай тул флаш карт үүсгэх, өөрийгөө сорих функц дутмаг байдаг.
%   \begin{figure}[H]
%     \centering
%     
\includegraphics[width=0.6\linewidth]{figures/chapter1/similarsystems/notionailogo.png}
%     \caption{Notion AI лого}
%     \label{fig:notionailogo}
%   \end{figure}
%   \item \textbf{Quizlet:} Флаш карт болон сорил үүсгэхэд хамгийн өргөн хэрэглэгддэг платформ. 2023 оноос "Q-Chat" хэмээх AI багшийг нэвтрүүлсэн. Гэсэн хэдий ч хэрэглэгч өөрөө контентоо оруулах шаардлагатай хэвээр байгаа бөгөөд урт хэмжээний PDF файлыг бүрэн задлан шинжлэхэд хязгаарлагдмал.
%   \begin{figure}[H]
%     \centering
%     \includegraphics[width=0.6\linewidth]{figures/chapter1/similarsystems/quizletlogo.png}
%     \caption{Quizlet лого}
%     \label{fig:quizletlogo}
%   \end{figure}
%   \item \textbf{Google NotebookLM:} 2024 онд танилцуулагдсан Google-ийн энэхүү хэрэгсэл нь эх сурвалжид суурилсан AI тэмдэглэлийн дэвтэр юм. Энэ нь хэрэглэгчийн оруулсан баримт бичиг дээр үндэслэн асуултад хариулж, хураангуй гаргадаг. Энэ нь манай системтэй хамгийн төстэй боловч сургалтын явцыг хянах болон шалгалт өгч сэдвийг "дууссан" гэж тэмдэглэх LMS-ийн шинж чанаргүй юм.
%   \begin{figure}[H]
%     \centering
%     \includegraphics[width=0.6\linewidth]{figures/chapter1/similarsystems/notebooklmlogo.png}
%     \caption{NotebookLM лого}
%     \label{fig:notebooklmlogo}
%   \end{figure}
%   \item \textbf{Mem.ai:} Хувийн мэдлэгийн сан үүсгэхэд чиглэсэн AI платформ. Мэдээллийг автоматаар холбох давуу талтай ч оюутны шалгалтад бэлдэх хэрэгцээг бүрэн хангадаггүй.
%   \begin{figure}[H]
%     \centering
%     
\includegraphics[width=0.6\linewidth]{figures/chapter1/similarsystems/memailogo.png}
%     \caption{Mem.ai лого}
%     \label{fig:memailogo}
%   \end{figure}
% \end{enumerate}

% \begin{longtable}{| p{0.28\textwidth} | p{0.12\textwidth} | p{0.12\textwidth} | p{0.14\textwidth} | p{0.12\textwidth} | p{0.14\textwidth} |}
% \caption{AI-д суурилсан сургалтын системүүдийн харьцуулалт}
% \label{tab:ai_learning_systems}\\
% \hline
% {\bfseries Үзүүлэлтүүд (Features)} & {\bf Notion AI} & {\bf Quizlet} & {\bf Google NotebookLM} & {\bf Mem.ai} & {\bf Proposed System} \\ \hline

% {\bf PDF Задлан шинжилгээ} & Хязгаарлагдмал & Үгүй & Тийм & Тийм & Тийм \\ \hline

% {\bf Автомат хураангуй (Summary)} & Тийм & Үгүй & Тийм & Тийм & Тийм \\ \hline

% {\bf Флаш карт \& Сорил үүсгэх} & Үгүй & Тийм (Гараар/AI) & Хязгаарлагдмал & Үгүй & Тийм (Автомат) \\ \hline

% {\bf Агуулгад суурилсан Чатбот} & Тийм (Ерөнхий) & Тийм (Q-Chat) & Тийм & Тийм & Тийм (Контексттэй) \\ \hline

% {\bf Суралцах явцын хяналт} & Үгүй & Тийм & Үгүй & Үгүй & Тийм \\ \hline

% {\bf Сургалтын бүтэц (Topic/Chapter)} & Үгүй & Үгүй & Үгүй & Үгүй & Тийм \\ \hline

% \end{longtable}

% \section{Технологийн судалгаа ба шийдэл}
% \paragraph{}

% Системийн хөгжүүлэлтийг хийхдээ үр ашигтай байдал болон кодын зохион байгуулалтыг сайжруулах үүднээс фронтэнд болон бакэнд хөгжүүлэлтийг нэгтгэсэн \textbf{Monorepo (Monolithic Repository)} архитектурыг ашиглахаар сонгосон. Энэ нь төслийн хувилбарын удирдлагыг хялбарчилж, хөгжүүлэлтийн хурдыг нэмэгдүүлэх давуу талтай.

% \subsection{Хиймэл оюун ухаан болон OCR технологи (AI \& OCR)}
% Гадаадын хиймэл оюун ухааны системүүд (жишээ нь: OpenAI, Claude) нь англи хэл дээр өндөр нарийвчлалтай ажилладаг боловч монгол хэлний дүрмийн онцлог, хэлц үгсийг бүрэн таних, ялангуяа монгол хэлээр бичигдсэн сканнердсан зурагнаас текст ялгах тал дээр ихээхэн алдаа гаргадаг. Иймд дотоодын онцлогт тохирсон дараах технологиудыг системд ашиглана:
% \begin{itemize}
%     \item \textbf{Egune LLM:} Монгол хэлний асар их өгөгдөл дээр сургагдсан үндэсний хиймэл оюун ухааны загвар. Оюутны оруулсан монгол хэл дээрх хичээлийн материалыг утгын алдаагүйгээр хураангуйлах, флаш карт болон сорилын асуулт үүсгэх гол цөм болно.
%     \begin{figure}[H]
%       \centering
%       
\includegraphics[width=0.6\linewidth]{figures/chapter1/egunelogo.png}
%       \caption{EguneAI лого}
%       \label{fig:egunelogo}
%     \end{figure}
%     \item \textbf{Chimege OCR:} Боловсролын салбарт түгээмэл ашиглагддаг сканнердсан ном, PDF болон зурган материалаас крилл монгол текстийг хамгийн өндөр нарийвчлалтайгаар таних хиймэл оюуны модел систем. Энэ нь хэрэглэгчдэд ямар ч форматын материалаас AI контент үүсгэх боломжийг олгоно.
%     \begin{figure}[H]
%       \centering
%       
\includegraphics[width=0.6\linewidth]{figures/chapter1/chimegelogo.png}
%       \caption{Chimege лого}
%       \label{fig:chimegelogo}
%     \end{figure}
% \end{itemize}

% \subsection{Backend хөгжүүлэлтэд ашиглах технологиуд}
% Систем нь нэгэн зэрэг олон хэрэглэгчийн PDF файлыг хүлээж авах, OCR ашиглан текст таних, AI модел руу хүсэлт илгээх зэрэг серверийн нөөц их шаардах үйлдлүүдийг гүйцэтгэнэ.
% \begin{itemize}
%     \item \textbf{Golang (Go хэл):} Маш өндөр хурдтай, санах ойн ашиглалт багатай бөгөөд зэрэгцээ ажиллагаа сайтай тул их хэмжээний файл боловсруулалтыг гацалтгүй хийхэд нэн тохиромжтой.
%     \item \textbf{Go-Chi Router:} Golang дээр суурилсан хөнгөн бөгөөд хурдтай замчлагч. REST API үүсгэх үйл явцыг хялбарчилна.
%     \item \textbf{WebSocket:} AI болон OCR боловсруулалт хийхэд тодорхой хугацаа шаардагдах тул боловсруулалтын ахиц хувийг хэрэглэгчийн дэлгэцэнд бодит хугацаанд дамжуулах зорилгоор ашиглана.
% \end{itemize}

% \subsection{Frontend хөгжүүлэлт}
% \begin{itemize}
%     \item \textbf{Vue.js болон Nuxt.js:} Вэб интерфэйсийг бүтээхэд зориулагдсан орчин үеийн фреймворк. Nuxt.js-ийг ашигласнаар хуудасны ачааллах хурд сайжирч, бүтэц зохион байгуулалт цэгцтэй болно.
%     \item \textbf{Shadcn-vue:} Хэрэглэгчийн интерфэйсийн (UI) бүрэлдэхүүн хэсгүүдийг хурдан бөгөөд загварлаг бүтээх боломжтой сан. Энэ нь responsive дизайн гаргах шаардлагыг бүрэн хангана.
% \end{itemize}

% \subsection{Өгөгдлийн сан болон Аюулгүй байдал}
% \begin{itemize}
%     \item \textbf{PostgreSQL:} Хэрэглэгчийн мэдээлэл, хичээлийн бүтэц, флаш карт, сорилын дүн зэрэг харилцан хамаарал бүхий өгөгдлүүдийг хадгалахад зориулагдсан, найдвартай реляцион өгөгдлийн сан.
%     \item \textbf{GORM:} Golang хэлнээс PostgreSQL-тэй харилцах харилцааг хялбарчлах Object-Relational Mapping (ORM) сан.
%     \item \textbf{JSON Web Token (JWT):} Хэрэглэгч болон Админы эрхийг ялгах, системд аюулгүй нэвтрэх үйл явцыг хангах зорилгоор ашиглана. 
% \end{itemize}

% \subsection{Дэд бүтэц болон Байршуулалт}
% \begin{itemize}
%     \item \textbf{Amazon Web Services (AWS):} Системийг интернэт орчинд байршуулж, тасралтгүй ажиллагааг хангах үүлэн технологийн үйлчилгээ.
%     \item \textbf{Nginx:} Сервер рүү ирэх урсгалыг хуваарилах (Reverse Proxy), HTTPS протоколын аюулгүй байдлыг хангах, мөн WebSocket-ийн холболтыг зохицуулах зорилгоор ашиглана.
% \end{itemize}

% \section{Бүлгийн дүгнэлт}
% \paragraph{}

% Нэгдүгээр бүлгийн хүрээнд уламжлалт сургалтын арга барилд тулгарч буй мэдээллийн хэт ачаалал, зурган (сканнердсан) материалын бэрхшээл болон өөрийгөө сорих хэрэглэгдэхүүний дутагдал зэрэг үндсэн асуудлуудыг онол, статистик судалгаанд үндэслэн тодорхойллоо. Эдгээр асуудлаас улбаалан системийг ашиглах үндсэн оролцогч талуудыг (оюутан, багш, админ) тодруулж, тэдний бодит хэрэгцээ шаардлагыг гаргаж ирсэн нь дараагийн бүлгүүдэд системийн шинжилгээ хийх бат бөх суурь болсон юм.
% Мөн олон улсын зах зээл дээрх Notion AI, Quizlet, Google NotebookLM зэрэг ижил төстэй системүүдэд харьцуулсан судалгаа хийхэд тэдгээр нь монгол хэл дээрх материал, ялангуяа крилл үсгээрх сканнердсан зурган өгөгдлийг танихдаа хязгаарлагдмал төдийгүй, сургалтын явцыг ухаалгаар хянах тогтолцоо дутмаг болох нь харагдсан. 

% Иймд энэхүү орон зайг нөхөх зорилгоор дотоодын Egune LLM болон Chimege OCR технологиудыг өндөр хурдны Golang (Backend), Vue.js/Nuxt.js (Frontend) зэрэг орчин үеийн шийдлүүдтэй хослуулан ашиглахаар сонголоо. Энэхүү технологийн бүрдэл нь монгол хэрэглэгчдийн онцлогт бүрэн нийцсэн, тасралтгүй бөгөөд хурдтай ажиллах ухаалаг сургалтын орчныг бий болгох бүрэн боломжтой юм. Дараагийн бүлэгт эдгээр онол, арга зүйн судалгаа, хэрэглэгчийн хэрэгцээнд тулгуурлан системийн нарийвчилсан шинжилгээ буюу үйл ажиллагааны болон үйл ажиллагааны бус шаардлагууд (FR, NFR), системийн архитектурыг дэлгэрэнгүй тодорхойлох болно.

\chapter{Системийн тухай онол, арга зүйн судалгаа} 

\section{Хичээлийн ухаалаг тэмдэглэлийн системийн тухай}
Орчин үеийн боловсролын салбарт мэдээллийн хэмжээ асар хурдацтай өсөж, оюутан сурагчид олон тооны ном, PDF материал, лекцийн тэмдэглэл зэргийг богино хугацаанд уншиж судлах шаардлагатай тулгарч байна. Уламжлалт аргаар тэмдэглэл хөтлөх, шалгалтад бэлдэх нь цаг хугацаа их шаарддагаас гадна сурлагын бүтээмжийг бууруулдаг. Иймд хиймэл оюун ухаан (AI) болон оптик тэмдэгт танигч (OCR) технологид суурилсан ухаалаг системийг хөгжүүлснээр сургалтын материалыг автоматаар боловсруулж, хураангуйлах, өөрийгөө сорих флаш карт (flashcard), асуулт хариулт үүсгэх замаар суралцах үйл явцыг илүү сонирхолтой, үр дүнтэй болгох боломжтой юм. Энэхүү систем нь зөвхөн тэмдэглэл хадгалах бус, суралцагчийн ахиц дэвшлийг хянах, сургалтын ухаалаг туслах байх зорилготой.

\section{Асуудлын тодорхойлолт}
Сургалтын үйл явц дахь уламжлалт арга барилд дараах нийтлэг асуудлууд тулгарч байна:
\begin{itemize}
    \item \textbf{Мэдээллийн хэт ачаалал:} Оюутнууд олон хуудас урт PDF материал уншиж, гол санааг ялгаж авахад маш их цаг зарцуулдаг. Боловсролын салбарын судалгаанаас үзэхэд оюутнуудын 70 гаруй хувь нь цагийн менежмент болон мэдээллийн хэт их багтаамжаас шалтгаалан өгөгдсөн унших материалыг бүрэн судалж амждаггүй бөгөөд энэ нь танин мэдэхүйн хэт ачаалалд (cognitive overload) хүргэдэг байна \cite{sweller2011cognitive, berry2010student}.
    
    \item \textbf{Сканнердсан болон зурган материалын бэрхшээл:} Багш нарын өгсөн материал эсвэл ном нь текст хэлбэрээр бус зураг байдлаар байх тохиолдол их бөгөөд үүнээс текст хуулж авах, хайлт хийх боломжгүй байдаг. Мэдээллийн технологийн салбарын тайланд дурдсанаар байгууллага болон боловсролын орчин дахь нийт дижитал мэдээллийн 80 орчим хувь нь бүтэцгүй өгөгдөл буюу зураг, сканнердсан баримт хэлбэрээр оршдог байна \cite{shukla2022unstructured}. Иймд OCR технологи ашиглахгүйгээр эдгээр материалыг бүрэн ашиглах боломж хязгаарлагдмал юм.
    
    \item \textbf{Өөрийгөө сорих хэрэглэгдэхүүн дутмаг:} Уншсан зүйлээ бататгахын тулд флаш карт болон сорилын асуулт гараар бэлтгэх нь оюутнаас давхар хөдөлмөр шаарддаг. Судлаач Ж.Данлоски (J. Dunlosky) нарын өргөн цар хүрээтэй судалгаагаар "өөрийгөө сорих" болон "завсарлагатай давтамж" нь хамгийн үр дүнтэй сургалтын арга болох нь батлагдсан боловч, оюутнуудын дийлэнх нь цаг хугацаа шаарддаг тул эдгээрээс зайлсхийж, зөвхөн материалыг идэвхгүй дахин унших аргыг сонгодог байна \cite{dunlosky2013improving}.
    
    \item \textbf{Хяналт ба тогтолцоогүй байдал:} Оюутнууд аль хичээлийн, аль сэдвийг бүрэн ойлгосон, юун дээр алдаж байгаагаа бодитоор хянаж чаддаггүй. Өөрийн суралцах явцыг хянадаг, системтэйгээр төлөвлөдөг оюутнуудын сурлагын амжилт бусдаас 20-30\%-иар өндөр байдаг хэдий ч, уламжлалт тэмдэглэлийн дэвтэр болон энгийн уншигч програмууд нь үүнийг автоматаар хянах боломжийг олгодоггүй \cite{zimmerman2002becoming}.
\end{itemize}

Эдгээр асуудлууд нь суралцагчийн цагийг үр дүнгүй үрж, сурах идэвхийг бууруулж байгаа тул AI болон OCR технологид суурилан эдгээрийг автоматжуулж шийдвэрлэх нэгдсэн систем зайлшгүй шаардлагатай байна.

\section{Оролцогч талууд}
Дээрх асуудлуудыг шийдвэрлэх системийг ашиглах, түүнд нөлөөлөх оролцогч талуудыг дараах байдлаар тодорхойлов:
\begin{enumerate}
    \item \textbf{Оюутан, суралцагч:} Системийн гол хэрэглэгч. Хичээлийн материалаа (текст болон зураг/OCR) оруулж, хураангуйлуулж, өөрийгөө сорьж шалгалт өгөх, чатботтой харилцах замаар сурлагын амжилтаа ахиулах зорилготой.
    \item \textbf{Багш, мэдлэг хуваалцагч:} Өөрийн бэлтгэсэн ухаалаг хичээл, тэмдэглэлийг "Public" төлөвт шилжүүлж, бусад оюутнуудад линкээр түгээх хүсэлтэй хэрэглэгч.
    \item \textbf{Системийн админ:} Системийн хэвийн ажиллагааг хангах, хэрэглэгчдийн дата болон AI боловсруулалтын хэрэглээ, статистикийг хянах, удирдах үүрэгтэй этгээд.
\end{enumerate}

\section{Ижил төстэй системийн судалгаа}
Дэлхийн зах зээл дээр хиймэл оюун ухаанд суурилсан тэмдэглэл хөтлөх болон сургалтын хэд хэдэн системүүд байдаг. Эдгээрийг судалж үзэх нь өөрийн системийн давуу талыг тодорхойлоход чухал юм.

\begin{enumerate}
  \item \textbf{Notion AI:} Notion нь тэмдэглэл хөтлөх хүчирхэг хэрэгсэл бөгөөд сүүлийн үед AI нэмэлтийг ашиглан текстийг хураангуйлах, засварлах боломжтой болсон. Гэвч энэ нь сургалтын үйл явцад зориулагдаагүй, ерөнхий зориулалттай тул флаш карт үүсгэх, өөрийгөө сорих функц дутмаг байдаг.
  \begin{figure}[H]
    \centering
    
\includegraphics[width=0.6\linewidth]{figures/chapter1/similarsystems/notionailogo.png}
    \caption{Notion AI лого}
    \label{fig:notionailogo}
  \end{figure}
  \item \textbf{Quizlet:} Флаш карт болон сорил үүсгэхэд хамгийн өргөн хэрэглэгддэг платформ. 2023 оноос "Q-Chat" хэмээх AI багшийг нэвтрүүлсэн. Гэсэн хэдий ч хэрэглэгч өөрөө контентоо оруулах шаардлагатай хэвээр байгаа бөгөөд урт хэмжээний PDF файлыг бүрэн задлан шинжлэхэд хязгаарлагдмал.
  \begin{figure}[H]
    \centering
    \includegraphics[width=0.6\linewidth]{figures/chapter1/similarsystems/quizletlogo.png}
    \caption{Quizlet лого}
    \label{fig:quizletlogo}
  \end{figure}
  \item \textbf{Google NotebookLM:} 2024 онд танилцуулагдсан Google-ийн энэхүү хэрэгсэл нь эх сурвалжид суурилсан AI тэмдэглэлийн дэвтэр юм. Энэ нь хэрэглэгчийн оруулсан баримт бичиг дээр үндэслэн асуултад хариулж, хураангуй гаргадаг. Энэ нь манай системтэй хамгийн төстэй боловч сургалтын явцыг хянах болон шалгалт өгч сэдвийг "дууссан" гэж тэмдэглэх LMS-ийн шинж чанаргүй юм.
  \begin{figure}[H]
    \centering
    \includegraphics[width=0.6\linewidth]{figures/chapter1/similarsystems/notebooklmlogo.png}
    \caption{NotebookLM лого}
    \label{fig:notebooklmlogo}
  \end{figure}
  \item \textbf{Mem.ai:} Хувийн мэдлэгийн сан үүсгэхэд чиглэсэн AI платформ. Мэдээллийг автоматаар холбох давуу талтай ч оюутны шалгалтад бэлдэх хэрэгцээг бүрэн хангадаггүй.
  \begin{figure}[H]
    \centering
    
\includegraphics[width=0.6\linewidth]{figures/chapter1/similarsystems/memailogo.png}
    \caption{Mem.ai лого}
    \label{fig:memailogo}
  \end{figure}
\end{enumerate}

\begin{longtable}{| p{0.28\textwidth} | p{0.12\textwidth} | p{0.12\textwidth} | p{0.14\textwidth} | p{0.12\textwidth} | p{0.14\textwidth} |}
\caption{AI-д суурилсан сургалтын системүүдийн харьцуулалт}
\label{tab:ai_learning_systems}\\
\hline
{\bfseries Үзүүлэлтүүд} & {\bf Notion AI} & {\bf Quizlet} & {\bf Google NotebookLM} & {\bf Mem.ai} & {\bf Бидний систем} \\ \hline

{\bf PDF Задлан шинжилгээ} & Хязгаарлагдмал & Үгүй & Тийм & Тийм & Тийм \\ \hline

{\bf Автомат хураангуй (Summary)} & Тийм & Үгүй & Тийм & Тийм & Тийм \\ \hline

{\bf Флаш карт \& Сорил үүсгэх} & Үгүй & Тийм (Гараар/AI) & Хязгаарлагдмал & Үгүй & Тийм (Автомат) \\ \hline

{\bf Агуулгад суурилсан Чатбот} & Тийм (Ерөнхий) & Тийм (Q-Chat) & Тийм & Тийм & Тийм (Контексттэй) \\ \hline

{\bf Суралцах явцын хяналт} & Үгүй & Тийм & Үгүй & Үгүй & Тийм \\ \hline

{\bf Сургалтын бүтэц (Topic/Chapter)} & Үгүй & Үгүй & Үгүй & Үгүй & Тийм \\ \hline
\end{longtable}

\section{Хэрэглэгчийн судалгаа}
Оролцогч талуудын хэрэгцээнд үндэслэн систем дараах үндсэн шаардлагуудыг хангасан байх ёстой. (Эдгээр нь 2-р бүлэгт заасан системийн FR, NFR шаардлагуудын үндэс болно).
\begin{itemize}
    \item \textbf{Суралцагчийн шаардлага:} 
    \begin{itemize}
        \item Нарийн төвөгтэй бүтэцтэй болон зураг агуулсан (OCR шаардлагатай) файлуудыг саадгүй боловсруулдаг байх.
        \item AI-аас хичээлийн агуулгын хураангуй, флаш карт, сорилын асуулт гаргуулж авдаг байх.
        \item Ойлгохгүй зүйлээ тухайн сэдэвтэй холбогдсон хиймэл оюун ухаант туслахаас асуудаг байх.
        \item Өөрийн суралцах явц буюу ахиц дэвшлийг системээс хянадаг, шалгалтад тэнцсэнээр хичээлийг үзэж дуусгасанд тооцдог байх.
    \end{itemize}
    \item \textbf{Багш, контент бүтээгчийн шаардлага:} Өөрийн үүсгэсэн материалыг олон нийтэд нээлттэй болгож хуваалцах боломжтой байх.
    \item \textbf{Админы шаардлага:} Системийн ачаалал, алдаа болон нийт хэрэглэгчдийн статистикийг харах, хяналтын самбартай байх.
\end{itemize}

\section{Технологийн судалгаа ба шийдэл}
Системийн хөгжүүлэлтийг хийхдээ үр ашигтай байдал болон кодын зохион байгуулалтыг сайжруулах үүднээс фронтэнд болон бакэнд хөгжүүлэлтийг нэгтгэсэн \textbf{Monorepo (Monolithic Repository)} архитектурыг ашиглахаар сонгосон. Энэ нь төслийн хувилбарын удирдлагыг хялбарчилж, хөгжүүлэлтийн хурдыг нэмэгдүүлэх давуу талтай.

\subsection{Хиймэл оюун ухаан болон OCR технологи (AI \& OCR)}
Гадаадын хиймэл оюун ухааны системүүд (жишээ нь: OpenAI, Claude) нь англи хэл дээр өндөр нарийвчлалтай ажилладаг боловч монгол хэлний дүрмийн онцлог, хэлц үгсийг бүрэн таних, ялангуяа монгол хэлээр бичигдсэн сканнердсан зурагнаас текст ялгах тал дээр ихээхэн алдаа гаргадаг. Иймд дотоодын онцлогт тохирсон дараах технологиудыг системд ашиглана:
\begin{itemize}
    \item \textbf{Egune LLM:} Монгол хэлний асар их өгөгдөл дээр сургагдсан үндэсний хиймэл оюун ухааны загвар. Оюутны оруулсан монгол хэл дээрх хичээлийн материалыг утгын алдаагүйгээр хураангуйлах, флаш карт болон сорилын асуулт үүсгэх гол цөм болно.
    \begin{figure}[H]
      \centering
      
\includegraphics[width=0.6\linewidth]{figures/chapter1/egunelogo.png}
      \caption{EguneAI лого}
      \label{fig:egunelogo}
    \end{figure}
    \item \textbf{Chimege OCR:} Боловсролын салбарт түгээмэл ашиглагддаг сканнердсан ном, PDF болон зурган материалаас крилл монгол текстийг хамгийн өндөр нарийвчлалтайгаар таних хиймэл оюун ухааны модель систем. Энэ нь хэрэглэгчдэд ямар ч форматын материалаас AI контент үүсгэх боломжийг олгоно.
    \begin{figure}[H]
      \centering
      
\includegraphics[width=0.6\linewidth]{figures/chapter1/chimegelogo.png}
      \caption{Chimege лого}
      \label{fig:chimegelogo}
    \end{figure}
\end{itemize}

\subsection{Backend хөгжүүлэлт}
Систем нь нэгэн зэрэг олон хэрэглэгчийн PDF файлыг хүлээж авах, OCR ашиглан текст таних, AI модель руу хүсэлт илгээх зэрэг серверийн нөөц их шаардах үйлдлүүдийг гүйцэтгэнэ.
\begin{itemize}
    \item \textbf{Golang (Go хэл):} Маш өндөр хурдтай, санах ойн ашиглалт багатай бөгөөд зэрэгцээ ажиллагаа сайтай тул их хэмжээний файл боловсруулалтыг гацалтгүй хийхэд нэн тохиромжтой.
    \item \textbf{Go-Chi Router:} Golang дээр суурилсан хөнгөн бөгөөд хурдтай замчлагч. REST API үүсгэх үйл явцыг хялбарчилна.
    \item \textbf{WebSocket:} AI болон OCR боловсруулалт хийхэд тодорхой хугацаа шаардагдах тул боловсруулалтын ахиц хувийг хэрэглэгчийн дэлгэцэнд бодит хугацаанд дамжуулах зорилгоор ашиглана.
\end{itemize}

\subsection{Frontend хөгжүүлэлт}
\begin{itemize}
    \item \textbf{Vue.js болон Nuxt.js:} Вэб интерфэйсийг бүтээхэд зориулагдсан орчин үеийн фреймворк. Nuxt.js-ийг ашигласнаар хуудасны ачааллах хурд сайжирч, бүтэц зохион байгуулалт цэгцтэй болно.
    \item \textbf{Shadcn-vue:} Хэрэглэгчийн интерфэйсийн (UI) бүрэлдэхүүн хэсгүүдийг хурдан бөгөөд загварлаг бүтээх боломжтой сан. Энэ нь responsive дизайн гаргах шаардлагыг бүрэн хангана.
\end{itemize}

\subsection{Өгөгдлийн сан болон Аюулгүй байдал}
\begin{itemize}
    \item \textbf{PostgreSQL:} Хэрэглэгчийн мэдээлэл, хичээлийн бүтэц, флаш карт, сорилын дүн зэрэг харилцан хамаарал бүхий өгөгдлүүдийг хадгалахад зориулагдсан, найдвартай реляцион өгөгдлийн сан.
    \item \textbf{GORM:} Golang хэлнээс PostgreSQL-тэй харилцах харилцааг хялбарчлах Object-Relational Mapping (ORM) сан.
    \item \textbf{JSON Web Token (JWT):} Хэрэглэгч болон Админы эрхийг ялгах, системд аюулгүй нэвтрэх үйл явцыг хангах зорилгоор ашиглана. 
\end{itemize}

\subsection{Дэд бүтэц болон Байршуулалт}
\begin{itemize}
    \item \textbf{Amazon Web Services (AWS):} Системийг интернэт орчинд байршуулж, тасралтгүй ажиллагааг хангах үүлэн технологийн үйлчилгээ.
    \item \textbf{Nginx:} Сервер рүү ирэх урсгалыг хуваарилах (Reverse Proxy), HTTPS протоколын аюулгүй байдлыг хангах, мөн WebSocket-ийн холболтыг зохицуулах зорилгоор ашиглана.
\end{itemize}

\section{Бүлгийн дүгнэлт}
Нэгдүгээр бүлгийн хүрээнд уламжлалт сургалтын арга барилд тулгарч буй мэдээллийн хэт ачаалал, зурган (сканнердсан) материалын бэрхшээл болон өөрийгөө сорих хэрэглэгдэхүүний дутагдал зэрэг үндсэн асуудлуудыг онол, статистик судалгаанд үндэслэн тодорхойллоо. Эдгээр асуудлаас улбаалан системийг ашиглах үндсэн оролцогч талуудыг (оюутан, багш, админ) тодруулж, тэдний бодит хэрэгцээ шаардлагыг гаргаж ирсэн нь дараагийн бүлгүүдэд системийн шинжилгээ хийх бат бөх суурь болсон юм.

Мөн олон улсын зах зээл дээрх Notion AI, Quizlet, Google NotebookLM зэрэг ижил төстэй системүүдэд харьцуулсан судалгаа хийхэд тэдгээр нь монгол хэл дээрх материал, ялангуяа крилл үсгээрх сканнердсан зурган өгөгдлийг танихдаа хязгаарлагдмал төдийгүй, сургалтын явцыг ухаалгаар хянах тогтолцоо дутмаг болох нь харагдсан. 

Иймд энэхүү орон зайг нөхөх зорилгоор дотоодын Egune LLM болон Chimege OCR технологиудыг өндөр хурдны Golang (Backend), Vue.js/Nuxt.js (Frontend) зэрэг орчин үеийн шийдлүүдтэй хослуулан ашиглахаар сонголоо. Энэхүү технологийн бүрдэл нь монгол хэрэглэгчдийн онцлогт бүрэн нийцсэн, тасралтгүй бөгөөд хурдтай ажиллах ухаалаг сургалтын орчныг бий болгох бүрэн боломжтой юм. Дараагийн бүлэгт эдгээр онол, арга зүйн судалгаа, хэрэглэгчийн хэрэгцээнд тулгуурлан системийн нарийвчилсан шинжилгээ буюу үйл ажиллагааны болон үйл ажиллагааны бус шаардлагууд (FR, NFR), системийн архитектурыг дэлгэрэнгүй тодорхойлох болно.